\subsection{A weight bound in each fixed cotangent homological degree}

\begin{proposition}[A-L10: Proved; fixed-degree prime-local support]
\label{AL10}
Put $S=\Gamma_{\mathbb Z}(T)$ and $R_N=S/I_N$, where
$N\geq2$ and $I_N=\bigoplus_{w\geq N}S^w$.
For a prime $p$ and an integer $n\geq0$, define
\[
 b_p(n)=\max\left(\{n+1\}\cup
 \left\{(n+2-2r)p^r:
                  1\leq r\leq\left\lfloor\frac{n+1}{2}\right\rfloor
 \right\}\right).
 \tag{AL.51}
\]
The second set is empty when $n=0$.
For every cellular free simplicial PD $S$-algebra
resolution $Q\longrightarrow R_N$, its special differential
module satisfies
\[
 H_n\!\left(N\Omega^1_{\mathrm{PD}}(Q/S)\right)^w
                  \otimes_{\mathbb Z}\mathbb Z_{(p)}=0
 \qquad\bigl(w>(N-1)(1+b_p(n))\bigr).
 \tag{AL.52}
\]
This includes $p=2$.  The bound depends on the fixed
homological degree, prime, and coefficient truncation.
\end{proposition}

\begin{proof}
Use first the actual finite-weight cotriple model
of A-S09-12.  Set
\[
 M_N=R_N^+=\bigoplus_{j=1}^{N-1}\mathbb Zt_j,
 \qquad G=\Gamma^+_{\mathbb Z},\qquad
 K_N=N\operatorname{Bar}_{\bullet}(\mathrm{id},G,M_N).
\]
The last bar face uses the PD multiplication of $R_N$;
the first face is the projection to arity one.
Filter this explicit bar complex by the number of
bottom letters.  Substitution preserves this number,
and a nonlinear application of the last face lowers
it.  Therefore the associated graded complex is the
bar construction on the same weighted group $M_N$
with zero nonlinear operations.  In each fixed
weight this filtration and its normalized chain
complex are finite, as proved by the nondegenerate
tree description in A-S09-12.

Here is the numerical specialization of the opened
Fu--Nuiten basis theorem needed for the bound
\cite[Theorem C, p.~6; Theorem 4.2, p.~40]{FuNuiten2026_check}.
A Lyndon word in the degree-zero letters has length
$L\geq1$, original weight $b$, and Lie degree $1-L$;
in particular $b\leq(N-1)L$.
Let $e$ and an operation sequence $(i_1,\ldots,i_r)$
label a basis class, and put $s=(1+e)L$.
Its bar degree and original weight are
\[
        n=s-1-\sum_{j=1}^r i_j+r,
        \qquad w=(1+e)p^r b\leq(N-1)s p^r.
 \tag{AL.53}
\]
When $r=0$, this gives $s=n+1$.
When $r>0$, the upper-entry family of that theorem
is empty: its terminal upper bound is at most zero,
whereas that family requires an entry greater than
one.  In the remaining family all $i_j<1$, hence
$i_j\leq0$, and $p i_{j+1}<i_j$.
An entry zero can occur only in the first position,
because a predecessor of zero would have to be
strictly positive.  Consequently
\[
        \sum_{j=1}^r i_j\leq-(r-1),\qquad
        n\geq s+2r-2,\qquad
        1\leq r\leq\left\lfloor\frac{n+1}{2}\right\rfloor.
 \tag{AL.54}
\]
These inequalities use no odd-prime restriction;
at $p=2$ the theorem sets $e=0$ and the same proof
applies.  Equations (AL.53)--(AL.54) prove that
the associated graded bar homology in degree $n$
has weight at most $(N-1)b_p(n)$.
The finite bottom-letter filtration therefore gives
\[
 H_n(K_N\otimes\mathbb F_p)^w=0
              \qquad(w>(N-1)b_p(n)).               \tag{AL.55}
\]
The weightwise bar duality used here is the explicit
finite-dimensional comparison (S09.18), applied to
the letters of weights $1,\ldots,N-1$; thus (AL.53)
computes this particular bar complex.

Now put
\[
 C_N=N\bigl(R_N\otimes_Q
                         \Omega^1_{\mathrm{PD}}(Q/S)\bigr).
\]
For $w>N$, the constant quotient $R_N\,dt_1$ has
no weight-$w$ part.  The coefficient-weight filtration
of $C_N^w$, constructed directly from special
differentiation in A-S09-12, has associated quotients
$K_N^{w-a}$ for $0\leq a<N$.  Its first face either
retains the coefficient weight or raises it; all
other faces preserve it.  If
$w>(N-1)(1+b_p(n))$, then every one of these
quotients has zero mod-$p$ homology in degree $n$
by (AL.55).  Finite successive long exact sequences
give
\[
             H_n(C_N^w\otimes\mathbb F_p)=0.       \tag{AL.56}
\]
The stated inequality also implies $w>N$, since
$(N-1)(1+b_p(n))\geq2(N-1)\geq N$.

For the chosen cotriple model, $C_N^w$ is a bounded
finite free integral complex.  Universal coefficients
inject $H_n(C_N^w)\otimes\mathbb F_p$ into the
zero group in (AL.56).  The finitely generated
$\mathbb Z_{(p)}$-module
$H_n(C_N^w)\otimes\mathbb Z_{(p)}$ thus has zero
quotient by $p$, so it is zero by Nakayama's lemma.
This proves the bound on the finite cotriple model.

For completeness, the comparison used to transfer
this conclusion is the actual one in A-S09-12.
The cellular module filtration gives a weak
equivalence
\[
 \Omega^1_{\mathrm{PD}}(Q/S)\longrightarrow
               R_N\otimes_Q\Omega^1_{\mathrm{PD}}(Q/S),
\]
whose cell quotients are the maps
$Q\otimes\mathbb Z[\Delta^j/\partial\Delta^j]
\to R_N\otimes\mathbb Z[\Delta^j/\partial\Delta^j]$.
For another cellular resolution $Q'$, the acyclic
augmentation kernels give actual PD maps
$u:Q\to Q'$ and $v:Q'\to Q$ over $S$ and $R_N$
and simplicial homotopies of their composites to
the identities.  Special differentiation followed
by $R_N$-base change preserves those homotopies.
Thus the two base-changed differential modules are
simplicially homotopy equivalent.  These maps prove
(AL.52) for every stated $Q'$ without asserting
that an arbitrary such resolution is finite in
each weight.
\end{proof}

The opposite quantifier pattern in A-S09-11 and
A-S09-12 is compatible with (AL.52): their nonzero
classes have both growing homological degree and
growing weight.  Equations (AL.53)--(AL.56) prove
the fixed-degree vanishing directly, rather than
inferring it from those examples.

\paragraph{Precise source locator.}
The opened primary version is
\url{https://arxiv.org/pdf/2610.36942v1}.
In its retained source
\texttt{Lie-October2026-ArXiv.tex},
Theorem C is at lines 425--440, the completely
unadmissible sequence definition at lines
1515--1529, and the weighted formula of Theorem
4.2 at lines 1757--1767.  The displayed inequalities
and the explicit bound (AL.51) are derived above.
